\section{Approach}



\subsection{Federated Model Formation}
In FL, multiple clients collaboratively train a global LLM ($G$) without sharing their raw data. Consider a FL round $r$ where $K$ clients participate. Each client $i \in \{1, \ldots, K\}$ possesses a local dataset $\mathcal{D}_i$ and performs local training on the global model from the previous round, producing an updated client model $C_i^{(r)}$. These client models are then aggregated by the central server to form the new global model $G^{(r)}$ for round $r$. Different model aggregation are available such as FedAvg and Fedprox. For instance FedAvg computes a weighted average of client model parameters:
\begin{equation}
G^{(r)} = \sum_{i=1}^K \rho_i^{(r)} \cdot C_i^{(r)}
\end{equation}
where $\rho_i^{(r)} = \frac{|\mathcal{D}_i|}{\sum_{j=1}^K |\mathcal{D}_j|}$ represents the aggregation coefficient proportional to client $i$'s dataset size.





\subsection{LLM Token Generation}
Once a global LLM is train it can deploy where it can take an input prompt $X$ to generate respnse $Y$. 


This prompt is first \textbf{tokenized} into a sequence $(x_1, \ldots, x_t)$, which is then mapped to continuous representations by a token embedding layer and fed into $G$. The core of the model is a \textbf{stack of $L$ identical Transformer Blocks}, applied sequentially. Each block contains trainable parameters and is composed of two primary sub-layers: a \textbf{multi-head self-attention} mechanism and a feed-forward network (MLP).

The self-attention mechanism enables the model to dynamically weigh the importance of different tokens relative to each other, allowing it to aggregate information from relevant positions to capture complex contextual dependencies. This is performed using multiple "heads" to learn different types of relationships simultaneously; a final \textbf{Output projection layer} then combines the information gathered by all heads. Following this, the \textbf{MLP} layer processes each token's representation individually. While self-attention handles contextual aggregation, the MLP provides a crucial non-linear transformation, allowing the model to store and apply its learned knowledge to refine the meaning of each token before it is passed to the next block.

After the full sequence is processed by $G$, the model produces logits (scores) for the next token over the entire vocabulary $V$. For autoregressive generation using greedy decoding, the next token $x_{t+1}$ is selected by choosing the token with the highest score:
% \begin{equation}
% x_{t+1} = \argmax_{v \in V} \left( G(x_1, \ldots, x_t) \right)_v
% \end{equation}

\begin{equation}
x_{t+1} = \operatorname*{arg\,max}_{v \in V} \left( G(x_1, \ldots, x_t) \right)_v
\end{equation}

This process continues iteratively. The generated token $x_{t+1}$ is appended to the input sequence, forming a new sequence $(x_1, \ldots, x_t, x_{t+1})$, which is then fed back into $G$ to generate the subsequent token $x_{t+2}$. This loop repeats until the model generates a special end-of-sequence token or reaches a maximum length.

\subsection{Problem Formulation}

For a given FL round $r$, we have access to the global model $G^{(r)}$ and the set of participating client models $\{C_1^{(r)}, C_2^{(r)}, \ldots, C_K^{(r)}\}$ that were retained before aggregation. When presented with an input prompt $\mathbf{x}$ and the corresponding response $\mathbf{y} = (t_1, t_2, \ldots, t_T)$ generated by $G^{(r)}$, our objective is to determine which client's training data contributed most significantly to the generation of $\mathbf{y}$. 


\subsection{Challenges} 
Developing an effective provenance method for federated LLMs requires addressing three interconnected challenges that arise from the unique characteristics of autoregressive language generation and transformer architectures.
This is extremely challenging as compared to a classification task. In classification, a model predicts a concrete output. For instance, in 10-class classification, the model only predicts one of the 10 classes. However, in the case of LLMs, the model response can contain thousands of tokens, and the response length will vary from one prompt to another. The attribution challenge is particularly acute for LLMs due to their extensive pre-training knowledge. When the global model generates token $y_t$, the output may stem from three potential sources: (1) knowledge acquired from client $i$'s federated training, (2) the model's pre-existing knowledge from pre-training, or (3) a combination thereof.Another challenge is that LLM parameters can range from millions to billions, which are hardly seen in traditional CNN architectures. Thus, enabling provenance at this scale of parameters is itself a challenging problem.





\textbf{Challenge 1: Variable-Length Autoregressive Generation.}
Unlike classification tasks that produce a single prediction, LLMs generate variable-length sequences through autoregressive token-by-token generation. Each token in the response may draw upon different aspects of the model's knowledge, potentially originating from different client training data. This raises a fundamental question: how do we attribute an entire response when different portions may come from different sources?

Our solution is to compute per-token provenance scores $\mathcal{P}_t(i)$ for each token $t_j \in \mathbf{y}$, which we then aggregate to obtain sequence-level attribution. The key insight is that the linear decomposition from FedAvg extends to each token generation step. For token $t_j$, the global model's computation decomposes as $\mathbf{o}_G^{(j)} = \sum_i \rho_i^{(r)} \cdot \mathbf{o}_{C_i}^{(j)}$, where each client's contribution can be measured independently. We aggregate these per-token contributions through summation: $\mathcal{P}(i) = \sum_{j=1}^T \mathcal{P}_{t_j}(i)$.



\textbf{Challenge 2: Knowledge Localization in Transformer Architectures.}
Given the linear decomposition from FedAvg, we could in principle measure client contributions at every neuron in every layer of the model. However, this would be computationally prohibitive and, more fundamentally, would conflate relevant and irrelevant contributions. Not all layers in a transformer LLM contribute equally to client-specific knowledge.

Modern transformer architectures organize computation hierarchically across layers, with different layers encoding different types of information. Research on neural network interpretability has revealed that early layers (bottom third) capture low-level linguistic features such as syntax, grammar, and basic semantic patterns. The final layers contains more fine grained knowledge. This hierarchical organization has direct implications for provenance, measuring provenance in late layers provides the strongest signal for attribution.

Within each transformer block, the architecture further separates functionality between attention mechanisms and feed-forward MLP sub-layers. Recent work has demonstrated that MLP blocks function as key-value memory stores that encode factual knowledge and associations learned during training. The self-attention mechanisms, in contrast, primarily perform information routing and contextual aggregation, with less direct storage of factual knowledge. This distinction suggests that MLP outputs in late layers should be prioritized for provenance measurement, as they directly represent the stored knowledge that differentiates clients.

Based on these insights, we focus provenance computation on the final two to three transformer blocks, specifically targeting MLP projection layers (gate, up, and down projections), the attention output projection layer, and the language modeling head. For a model with $L$ layers, we select layers from the set $\{\ell : \ell \geq L - 3\}$. This selective approach reduces computational overhead by over 80\% compared to measuring all layers, while concentrating on the components that carry the strongest client-specific signals. The attention output projection is included because, while attention performs routing, the output projection integrates multi-head results and can exhibit client-specific patterns when clients have distinct context-handling styles. The language modeling head is critical as it directly determines token predictions and carries strong provenance signals as the final transformation before vocabulary logits.

\textbf{Challenge 3 Relevance Filtering Through Gradient Weighting.}
Even when focusing on late-layer neurons, a fundamental challenge remains: not all neurons are relevant for generating a particular token. A client might have large activations in neurons that are completely irrelevant to the current token prediction. 

The naive approach of directly computing $\mathbf{o}_{C_i} = \theta_{C_i}^{(r)} \cdot \mathbf{x}$ for each client would conflate relevant and irrelevant activations, producing noisy attribution. We need a mechanism to automatically identify which neurons matter for the specific token being generated, and weight client contributions accordingly.

Our solution leverages the gradient of the token prediction with respect to layer activations as an importance weighting mechanism. For a specific layer $\ell$ and token $t_j$, we compute the gradient $\mathbf{g}^\ell = \frac{\partial \log p(t_j | \mathbf{x}_{<j}; G^{(r)})}{\partial \mathbf{a}_G^\ell}$, where $\mathbf{a}_G^\ell$ is the activation of layer $\ell$ in the global model. This gradient quantifies how much each dimension of the layer's activation influences the final token prediction. Dimensions with larger gradient magnitudes are more influential in determining the output, while dimensions with zero or near-zero gradients are irrelevant regardless of their activation magnitude.

The key insight is that the inner product $\langle \mathbf{a}_{C_i}^\ell, \mathbf{g}^\ell \rangle$ automatically performs relevance-weighted attribution. For each dimension $k$ in the activation space, the contribution is $a_{C_i,k}^\ell \cdot g_k^\ell$. When gradient $g_k^\ell$ is zero, the contribution is zero regardless of how large $a_{C_i,k}^\ell$ might be. Conversely, when $g_k^\ell$ is large, the client's activation at that dimension directly influences the provenance score. This selective weighting focuses attribution on task-relevant computations, filtering out irrelevant neurons that would otherwise inject noise into the attribution.

The mathematical intuition is that $\mathbf{g}^\ell$ defines an "importance direction" in the activation space—the direction most responsible for producing token $t_j$. The client whose activation $\mathbf{a}_{C_i}^\ell$ aligns most with this direction is the one whose knowledge most influenced that token generation. This gradient-based weighting is why our approach produces more accurate provenance than methods that rely solely on activation magnitudes or parameter similarities.

\subsection{Insights}
To address this problem, we propose the idea of \emph{token provenance}. The idea is we track the provenance of each generated token, and the client who has the maximum contribution to the whole generated response is the one responsible for that response. Our goal is to produce an attribution vector $\mathcal{P} \in \mathbb{R}^K$ where $\mathcal{P}(i)$ quantifies the contribution of client $i$ to the generation of response $\mathbf{y}$. To enable fine-grained analysis, we compute per-token provenance scores and aggregate them to obtain sequence-level attribution. The client with the highest attribution score is identified as the primary source of the generated response.However, enabling provenance at each weight across billions of parameters quickly explodes the problem. Tracking provenance for every parameter across $K$ clients—even a small number like 5 or 10—and repeating this computation for each generated token is computationally intractable.To make this problem tractable, we leverage key structural insights from the Transformer architecture. We carefully selected specific layers for provenance tracking based on their functional roles, rather than monitoring all parameters.

\begin{itemize}
    \item \textbf{Self-Attention Sub-layer:} The self-attention mechanism, while containing many parameters in its query (Q), key (K), and value (V) projections, ultimately consolidates all information from its multiple heads into the Output Projection Layer. This layer acts as a bottleneck that merges and refines the complete contextual knowledge aggregated by the attention mechanism. We hypothesize that tracking provenance at this single layer is sufficient to capture the attention block's contribution, thus avoiding the massive overhead of tracking individual Q, K, and V computations.
    
    \item \textbf{MLP Sub-layer:} Similarly, the MLP (feed-forward network) block, which stores and applies factual knowledge, consists of multiple linear layers. We posit that the final layer in the MLP sequence contains the richest, most refined representation before the output is passed to the next Transformer block.
\end{itemize}

By restricting our provenance analysis to only these two critical layers (the Output Projection layer and the final MLP layer) within each Transformer block, we drastically reduce the number of parameters to track. This approach exploits the hierarchical and repetitive structure of Transformers to make token-level provenance in federated LLMs computationally feasible. In our evaluation, we will empirically validate this design, demonstrating that this efficient setting achieves high attribution accuracy while imposing minimal computational overhead. Next disscuss how to enable provenace across these layers and each token. 



\subsection{Federated Averaging and Linear Decomposition}

The Federated Averaging (FedAvg) algorithm, which serves as the aggregation mechanism in our framework, possesses a crucial mathematical property that enables provenance tracking. Understanding this property is essential to our approach.

At each round $r$, after local training, the server aggregates the client models using weighted averaging:
\begin{equation}
\theta_{\text{global}}^r = \sum_{i \in \mathcal{S}_r} w_i^r \cdot \theta_i^r
\end{equation}
where $w_i^r = \frac{|\mathcal{D}_i|}{\sum_{j \in \mathcal{S}_r} |\mathcal{D}_j|}$ represents the aggregation weight proportional to client $i$'s dataset size, and $\sum_{i \in \mathcal{S}_r} w_i^r = 1$.

The key insight is that this weighted averaging creates a \textit{linear composition} of client model parameters. Consider a single neuron or weight matrix $\theta$ within the model (suppressing round superscript for clarity). The global model's parameter is:
\begin{equation}
\theta_{\text{global}} = \sum_{i=1}^N c_i \cdot \theta_i
\end{equation}
where $c_i$ represents the cumulative contribution coefficient of client $i$ after multiple rounds of aggregation (with $c_i = 0$ for clients not yet participated).

This linear composition has profound implications for inference. When the global model processes an input $\mathbf{x}$, the computation at this neuron is:
\begin{equation}
\begin{split}
\mathbf{o}_{\text{global}} &= \theta_{\text{global}} \cdot \mathbf{x} \\
&= \left(\sum_{i=1}^N c_i \cdot \theta_i\right) \cdot \mathbf{x} \\
&= \sum_{i=1}^N c_i \cdot (\theta_i \cdot \mathbf{x}) \\
&= \sum_{i=1}^N c_i \cdot \mathbf{o}_i
\end{split}
\end{equation}
where $\mathbf{o}_i = \theta_i \cdot \mathbf{x}$ represents what client $i$'s model would compute independently on the same input.

This decomposition reveals that the global model's output is literally a weighted sum of individual client outputs. The global model's computation can be viewed as an ensemble of client-specific computations, linearly combined through the aggregation weights. This mathematical structure is what makes provenance tractable: by analyzing how much each client's hypothetical computation $\mathbf{o}_i$ contributes to the final output, we can attribute predictions back to their source clients.

However, directly computing $\mathbf{o}_i = \theta_i \cdot \mathbf{x}$ for each client requires maintaining all client models, which is impractical. Our gradient-based approach circumvents this by leveraging the activation patterns and gradient signals during inference to implicitly measure these contributions without explicitly storing all client parameters.

\subsection{Gradient-Based Neuron Provenance}

Our provenance method operates during the autoregressive token generation process. For each generated token, we compute provenance scores by analyzing the relationship between client-specific model activations and the gradient of the output with respect to those activations.

\subsubsection{Single Token Provenance}

Consider the generation of a single token $y_t$ at position $t$. Let $\mathbf{x}_{<t} = (y_1, \ldots, y_{t-1})$ denote the context (previous tokens). The global model computes the next token distribution as $p(y_t | \mathbf{x}_{<t}; \theta_{\text{global}})$, and selects $y_t = \arg\max p(y_t | \mathbf{x}_{<t}; \theta_{\text{global}})$.

For a specific layer $\ell$ in the model (e.g., an MLP block or attention output projection), let:
\begin{itemize}
\item $\mathbf{a}_{\text{global}}^\ell \in \mathbb{R}^d$ denote the activation (output) of layer $\ell$ when processing input $\mathbf{x}_{<t}$ through the global model
\item $\mathbf{g}^\ell = \frac{\partial \log p(y_t | \mathbf{x}_{<t}; \theta_{\text{global}})}{\partial \mathbf{a}_{\text{global}}^\ell} \in \mathbb{R}^d$ denote the gradient of the log-probability of the selected token $y_t$ with respect to the layer's activation
\end{itemize}

The gradient $\mathbf{g}^\ell$ quantifies how much each dimension of the layer's activation influences the final token prediction. Intuitively, dimensions with larger gradient magnitudes are more influential in determining the output.

Now, for each client $i$, we compute what that layer would output if the client's model weights were used instead of the global weights, while keeping the input activation from the global model's forward pass. Specifically, let $\mathbf{a}_i^\ell$ denote the output of layer $\ell$ using client $i$'s parameters $\theta_i^\ell$ on the global model's input to that layer. This represents the client-specific activation pattern for the same context.

The provenance score of client $i$ at layer $\ell$ for token $y_t$ is computed as the inner product between the client's activation and the output gradient:
\begin{equation}
S_i^{\ell}(t) = \langle \mathbf{a}_i^\ell, \mathbf{g}^\ell \rangle = \sum_{j=1}^d a_{i,j}^\ell \cdot g_j^\ell
\end{equation}

This inner product measures how aligned the client's activation is with the gradient direction. A large positive value indicates that client $i$'s parameters at layer $\ell$ push the model toward generating token $y_t$, while a negative or small value indicates little contribution.

The mathematical intuition is that $\mathbf{g}^\ell$ defines a "importance direction" in the activation space—the direction most responsible for producing $y_t$. The client whose activation $\mathbf{a}_i^\ell$ aligns most with this direction is the one whose knowledge most influenced that token generation.

\subsubsection{Multi-Layer Aggregation}

Modern LLMs consist of multiple layers, each contributing differently to the final output. To obtain a comprehensive provenance score, we aggregate contributions across multiple selected layers. Let $\mathcal{L}$ denote the set of layers chosen for provenance computation (we discuss layer selection in Section 2.4). The total provenance score for client $i$ on token $y_t$ is:
\begin{equation}
S_i(t) = \sum_{\ell \in \mathcal{L}} S_i^\ell(t) = \sum_{\ell \in \mathcal{L}} \langle \mathbf{a}_i^\ell, \mathbf{g}^\ell \rangle
\end{equation}

This summation accumulates evidence from different layers. Each layer captures different aspects of the model's computation—attention layers aggregate contextual information while MLP layers store factual knowledge. By aggregating across layers, we obtain a more robust provenance signal.

\subsubsection{Sequence-Level Provenance}

For a complete generated sequence $y = (y_1, \ldots, y_T)$, we aggregate the per-token provenance scores to obtain an overall attribution for the entire response:
\begin{equation}
S_i = \sum_{t=1}^T S_i(t) = \sum_{t=1}^T \sum_{\ell \in \mathcal{L}} \langle \mathbf{a}_i^\ell(t), \mathbf{g}^\ell(t) \rangle
\end{equation}

where $\mathbf{a}_i^\ell(t)$ and $\mathbf{g}^\ell(t)$ denote the activations and gradients at token position $t$.

Finally, we normalize these scores using softmax to obtain a probability distribution over clients:
\begin{equation}
P_i = \frac{\exp(S_i)}{\sum_{j=1}^N \exp(S_j)}
\end{equation}

The client with the highest probability $\hat{i} = \arg\max_i P_i$ is identified as the primary source of the generated response. This attribution provides explainability for federated model outputs and enables detection of malicious contributors in adversarial scenarios.